% ============================================================
% 第 19 章 同步电动机和同步补偿机（教材第五篇，印刷页 215--225）
% 节名：第一节 同步电动机电压方程和相量图；第二节 同步电动机的功角特性与 V 形曲线；
%       第三节 同步电动机的启动与调速；第四节 同步补偿机；第五节 永磁同步电动机
% ============================================================
\chapter{同步电动机和同步补偿机}

同步电机是可逆的：同一台机器，输入机械功率、输出电功率时为发电机；输入电功率、输出机械功率时为电动机。本章把第 17 章的方程改写成\textbf{电动机惯例}，重点解决两个新问题：\textbf{怎么起动}（同步电动机没有起动转矩）和\textbf{V 形曲线}（电流随励磁变化）。

\section{同步电动机电压方程和相量图}

\subsection{正方向的改变}
改作电动机运行时，电枢电流的正方向改为\textbf{从电网流入电机}，与感应电动势 $\dot E_0$ 反方向（吸收电功率），即
\begin{equation}
\dot U=\dot E_0+\dot I\left(R_a+\mathrm{j}X_t\right),
\end{equation}
或写成
\begin{equation}
\dot E_0=\dot U-\dot I\left(R_a+\mathrm{j}X_t\right).
\end{equation}

\keyformula{\dot U=\dot E_0+\dot I\left(R_a+\mathrm{j}X_t\right)\quad\text{（电动机惯例）}}

与发电机方程相比，只是 $\dot I$ 的符号变了。这意味着：\textbf{同一台电机，发电机与电动机的方程是同一个，区别只在电流方向与功率流向}。

\subsection{相量图与功率角}
电动机的功率角 $\delta$ 为 $\dot E_0$ 滞后 $\dot U$ 的角度（取负值或按绝对值定义），表示转子磁极轴线\textbf{滞后}于气隙磁场轴线。物理图像：电动机拖着负载，转子"落后"于旋转磁场，靠电磁"弹簧"拉着走。

电枢电流与端电压的夹角即功率因数角 $\varphi$；$\dot E_0$ 与 $\dot I$ 的夹角为 $\psi$。电动机可以工作在\textbf{超前}功率因数（过励）状态——这是同步电动机相对感应电动机的最大优点：可以补偿电网的无功。

\section{同步电动机的功角特性与 V 形曲线}

\subsection{功角特性}
电动机的电磁功率（由电网输入、转换为机械功率的部分）为
\begin{equation}
P_{em}=\frac{mE_0U}{X_t}\sin\delta ,
\end{equation}
凸极机为
\begin{equation}
P_{em}=\frac{mE_0U}{X_d}\sin\delta+\frac{mU^{2}}{2}\left(\frac{1}{X_q}-\frac{1}{X_d}\right)\sin2\delta .
\end{equation}
形式上与发电机完全相同，只是 $\delta$ 取相反符号（功率流向相反）。最大电磁功率（过载能力）同样出现在 $\delta=90^\circ$ 附近。

\subsection{V 形曲线}
在\textbf{输出功率恒定}（即输入有功不变）的条件下，改变励磁电流 $I_f$，电枢电流 $I$ 随之变化，其曲线形如字母"V"，故称 \textbf{V 形曲线}。

\subsubsection{形成机理}
输出功率不变意味着
\begin{equation}
P=\sqrt3UI\cos\varphi=\text{常数}\quad\Longrightarrow\quad I\cos\varphi=\text{常数},
\end{equation}
即电枢电流的\textbf{有功分量不变}。而励磁改变的是 $E_0$，从而改变\textbf{无功分量} $I\sin\varphi$。于是
\begin{equation}
I=\sqrt{\left(I\cos\varphi\right)^{2}+\left(I\sin\varphi\right)^{2}},
\end{equation}
当无功分量为零（$\cos\varphi=1$）时 $I$ 最小，即为 V 形曲线的\textbf{最低点}（正常励磁）；励磁大于该值时过励（电流超前、输出容性无功），小于该值时欠励（电流滞后、吸收感性无功）。

\subsubsection{不同负载下的 V 形曲线}
对应不同的输出功率可画出一族 V 形曲线：负载越大，曲线整体上移、最低点电流越大；各曲线最低点的连线称为\textbf{临界励磁线}（$\cos\varphi=1$ 线），其右侧为过励区、左侧为欠励区。欠励过度时，电机会因励磁不足而\textbf{失步}，故存在\textbf{稳定极限}（对应 $\delta=90^\circ$），V 形曲线左侧的部分区域是不稳定区。

\begin{center}
\begin{tabular}{@{}p{2.2cm}p{3.0cm}p{4.0cm}p{3.4cm}@{}}
\toprule
励磁状态 & 电枢电流性质 & 对电网的作用 & 典型用途 \\
\midrule
正常励磁 & $\cos\varphi=1$，$I$ 最小 & 不交换无功 & V 形曲线最低点 \\
过励 & 电流超前（容性） & 向电网输出感性无功 & 改善功率因数（进相运行） \\
欠励 & 电流滞后（感性） & 从电网吸收感性无功 & 一般不采用，易失步 \\
\bottomrule
\end{tabular}
\end{center}

\corepoint{V 形曲线的关键前提是"\textbf{有功不变}"。若输出功率变了，曲线整体平移，最低点位置也变——所以问"增大励磁电流，电流如何变化"时，必须先确认"输出功率是否不变"。另外记住：\textbf{同步电动机过励时功率因数超前}，这正是它能当"无功补偿装置"用的原因，也是它与感应电动机最本质的差别。}

\subsection{同步电动机的优缺点}
\begin{itemize}
  \item \textbf{优点}：转速严格恒定（$n=n_1$，与负载无关）；功率因数可调，可过励运行补偿电网无功；过载能力强（气隙大）、效率高。
  \item \textbf{缺点}：需要直流励磁电源；\textbf{没有起动转矩}，必须采用专门的起动方法；结构复杂、造价高。
\end{itemize}

\section{同步电动机的启动与调速}

\subsection{为什么不能自起动}
稳态运行时转子以同步转速旋转，转子磁极与气隙磁场相对静止。若在转子静止时投入三相电源，气隙旋转磁场以 $n_1$ 旋转，转子受到的电磁转矩在一个周期内\textbf{平均值为零}（时而正、时而负），转子无法加速，故同步电动机不能自起动。

\subsection{起动方法}
\begin{enumerate}
  \item \textbf{异步起动法}（最常用）：在转子上加装\textbf{笼型起动绕组}（阻尼绕组）。起动时先不加励磁，定子接三相电源，电机作为笼型感应电动机起动加速；当转速接近同步转速（约 95\% $n_1$）时投入励磁，靠同步转矩把转子\textbf{牵入同步}。牵入同步的瞬间会产生振荡，阻尼绕组可抑制振荡。
  \item \textbf{辅助电动机起动法}：用一台小电动机把同步电动机拖到接近同步转速，再投入电网与励磁。设备多、操作繁，仅用于特殊场合。
  \item \textbf{变频起动法}：用变频装置从低频开始供电，使旋转磁场转速从零缓慢上升，转子始终同步跟随。起动平滑、无冲击，是现代大容量同步电动机（如轧机、磨机）的主流方案。
\end{enumerate}

\subsection{调速}
同步电动机的转速 $n=60f/p$，只能通过\textbf{变频}（改变 $f$）或\textbf{变极}（改变 $p$）调速，且为有级（变极）或无级（变频）。变频调速时同样要保持 $U/f$ 恒定以维持磁通。这比感应电动机的调速手段少（因为没有转差率可调），但转速精度最高。

\section{同步补偿机}

同步补偿机（调相机）是\textbf{轴上不带任何机械负载}的同步电机，专门用来向电网输出或吸收无功功率。

\begin{itemize}
  \item \textbf{过励运行}：向电网输出感性无功，相当于并联电容器，可提高电网功率因数、支撑受端电压；
  \item \textbf{欠励运行}：从电网吸收感性无功，相当于并联电抗器，可抑制长线路空载或轻载时的工频过电压。
\end{itemize}
调相机的优点是无功调节连续、平滑、容量大；缺点是旋转设备有损耗、维护量大、造价高。近年来大量被\textbf{静止无功补偿器（SVC）}与\textbf{静止同步补偿器（STATCOM）}取代，但在特高压长距离输电中仍有应用。

\section{永磁同步电动机}

用\textbf{永磁体}代替励磁绕组建立主磁场，取消励磁电源、滑环与电刷，是同步电机的重要发展方向。

\subsubsection{结构特点}
按永磁体在转子上的位置，分为\textbf{表贴式}（永磁体贴在转子表面，$L_d\approx L_q$）与\textbf{内置式}（永磁体嵌入转子内部，$L_d<L_q$，存在磁阻转矩）。永磁材料常用钕铁硼（NdFeB）、钐钴等。

\subsubsection{优点}
\begin{itemize}
  \item 无励磁损耗，效率高（比同容量感应电动机高 2\%～8\%）；
  \item 功率密度大、体积小、质量轻；
  \item 无电刷滑环，可靠性高、维护简单；
  \item 配合变频器可实现宽范围调速与精确的位置/转矩控制。
\end{itemize}

\subsubsection{问题}
\begin{itemize}
  \item 永磁体价格高，且高温下可能\textbf{失磁}（需控制温度）；
  \item 无法通过调节励磁来调压、调无功，弱磁调速范围受限；
  \item 停机检修时反电动势仍然存在，安全上需注意。
\end{itemize}
永磁同步电动机广泛用于新能源汽车驱动、伺服系统、电梯、压缩机、风电（直驱永磁）等领域，是当前电机技术最活跃的方向之一。

\begin{center}
\begin{tabular}{@{}p{2.2cm}p{3.6cm}p{3.8cm}p{3.4cm}@{}}
\toprule
项目 & 同步电动机 & 感应电动机 & 直流电动机 \\
\midrule
转速 & 严格 $n_1$，与负载无关 & 随负载略降（有转差） & 随负载下降较多 \\
功率因数 & 可调（可超前） & 滞后，轻载时很低 & 较高（有换向器） \\
起动 & 需异步/变频起动 & 可自起动 & 可自起动，需限流 \\
调速 & 变频/变极 & 变频/变极/变转差 & 调压/调磁，最方便 \\
维护 & 需励磁（永磁机免） & 笼型免维护 & 换向器电刷需维护 \\
\bottomrule
\end{tabular}
\end{center}
